\section{Introduction}\label{sec:tsps-introduction}

Structure-preserving signature (SPS) schemes, in which messages, public keys and
signatures are all group elements, have proven useful across a range of
privacy-preserving applications, including blind signatures, whose signers do
not see what they sign, delegatable anonymous credentials, which holders can
pass on to others, and group signatures, which hide which member of a group
signed~\cite{AFRICACRYPT:FucVer10,CCS:CamDriDub17,RSA:CriLys19,PoPETS:MSBM23,AC:Groth06,AC:Groth07,C:LibPetYun15}.

These applications rely on signing group elements directly. A signer can sign a
commitment without opening it, and the holder of the signature can later prove
statements about the committed value in zero knowledge without revealing it. An
ordinary signature scheme hashes the message first, so proving knowledge of a
signature on a commitment needs a general-purpose proof over the hash function.
Signing the group elements directly keeps these proofs algebraic and efficient.

In systems with entities that publicly certify state, such as transparency log
providers or ledger administrators, threshold certification shares the signing
key between several signers, so that a threshold of them is needed to sign a
commitment and no single signer or compromised key share can certify alone.
Verifiers still see one public key and ordinary signatures.

A ledger records
commitments to tokens or accounts, and to spend a token or report on an account,
its owner needs a signature showing that the commitment is
recorded~\cite{EPRINT:ACDDET19,EPRINT:TBAAGP22,CCS:KiaKohSar22,RSA:PoiSan18,NDSS:SABMD19}.
With an SPS scheme, the certifier signs the committed group elements as they
appear on the ledger and publishes the signatures, without interacting with the
owner, and the owner later proves statements about the committed values in zero
knowledge. Where a ledger is run by several authorities, for example a central
bank digital currency in which every transaction is endorsed before it is
recorded, a threshold signature scheme divides the certification between them.

Certificate transparency (CT) and key
transparency (KT) logs are append-only, publicly verifiable records of
certificates or keys. Each entry is sent to a log provider, signed, and appended
to the log. In CT the log server returns a signed certificate timestamp to the
certificate provider, and in KT the log server signs the mapping from users'
identities to their public keys. Transparency logs move trust from the application
serving a key onto the log providers, and a threshold signature scheme narrows
what any one of them can do. Signatures from a threshold SPS scheme, given to the
owners of keys, would also let them prove, in zero knowledge, that their keys are
recorded in the log without revealing which keys they are.

\paragraph{Contributions.} We give two threshold SPS schemes, built from
existing SPS schemes and secure multiparty computation (MPC)
techniques~\cite{C:AGHO11,C:KilPanWee15,C:DPSZ12,CCS:KelOrsSch16,C:GLOPS21}.
One of our instantiations has signatures of three group elements, the smallest
possible for SPS in type III groups, but is secure only in the
GGM~\cite{C:AGHO11}. The other relies only on the SXDH
assumption~\cite{C:KilPanWee15}.

Both schemes start from a common construction, which applies to any SPS scheme
with decomposable signing, a property we define. All message-independent
computation runs under MPC in an offline phase, and once the message is known,
signing needs no interaction. \Cref{sec:tsps-security} reduces security to the
unforgeability of the underlying scheme.

Our second contribution is an implementation. Computation on shared group
elements has been implemented for two parties~\cite{NDSS:GPVRNC23,LC:ADEO21}
and for many parties on curves without a pairing~\cite{ESORICS:DOKSS20}. We add
sharings of pairing-group elements, and raising them to shared exponents, to
MP-SPDZ, a library implementing a wide range of MPC protocols, for any number of
parties~\cite{CCS:Keller20}. To our knowledge, this is
the first multiparty implementation of computation on shared elements of pairing
groups.

\subsection*{Related work}

\paragraph{SPS schemes.} The first SPS scheme signs messages of arbitrary length
with signatures of seven group elements, and is secure in the standard model
under a non-interactive assumption~\cite{C:AFGHO10}. Interactive assumptions,
which give the adversary an oracle, such as the assumption of Pointcheval and
Sanders, are avoided where possible, as they come close to assuming the security
of the scheme itself~\cite{RSA:PoiSan16}. A signature in type III groups
needs at least three group elements, and there are matching constructions in the
GGM, among them the rerandomisable scheme $\AGHO$ that we use~\cite{C:AGHO11}.
SPS schemes with such short signatures cannot be proven secure from non-interactive assumptions
by a large class of reductions~\cite{AC:AbeGroOhk11}. The scheme $\KPW$ has
signatures of seven group elements and is secure under SXDH~\cite{C:KilPanWee15}.
Fully structure-preserving
signature schemes also have secret keys that are group elements, so
a user can prove statements about their own key with Groth-Sahai proofs, which
prove pairing product equations in zero knowledge, but their signatures grow with the square root of the message
length~\cite{JC:AGKOT19,EC:GroSah08}.

\paragraph{Threshold SPS schemes.} Threshold SPS schemes have only recently
been considered. One non-interactive construction is restricted to signing
indexed Diffie-Hellman tuples, of the form $(\gen^\msg,
\ggen^\msg)$~\cite{AC:CKPSS23}. Existing commitments and public keys lie in one group, so it
cannot sign them, which rules out the ledger and transparency-log uses above.


Following our work, a non-interactive threshold SPS scheme was given that signs
arbitrary group elements and is secure under SXDH~\cite{PKC:MMSST24}. It builds
on $\KPW$, and its signatures consist of seven group elements, against three for
threshold $\AGHO$. We benchmark against it in \Cref{sec:tsps-evaluation}. Building on this scheme, a threshold SPS
scheme with silent setup lets signers generate their keys independently, with no
interactive key generation, at the cost of larger signatures, a trusted
setup, whose public parameters come from secrets that must be deleted, and security in the GGM~\cite{EPRINT:GGPR25}.

\paragraph{Layout.} The rest of the chapter is organised as follows. \Cref{sec:tsps-definitions} defines threshold SPS schemes
and the MPC interfaces our constructions use.
\Cref{sec:tsps-construction} gives the offline and online decomposition and the
generic construction, and \Cref{sec:tsps-instantiation} instantiates it with
$\AGHO$ and $\KPW$. \Cref{sec:tsps-security} proves both secure,
\Cref{sec:tsps-evaluation} evaluates them, and \Cref{sec:tsps-summary}
concludes.
